\subsection{R 中区间的性质}

命题 2. R 中每个闭（相应地开）区间是 R 中的闭（相应地开）集。

集合 \([a, \rightarrow[ = a + \mathbf{R}_+ \text{ 与 } ]\leftarrow, a] = a - \mathbf{R}_+\) 分别由 \(\mathbf{R}_+\) 与 \(-\mathbf{R}_+\) 经平移得到，因而是闭集（第三章 § 1 第 1 目）；集合 \(]\leftarrow, a[\) 与 \(]a, \rightarrow[\) 是它们的余集，因而是开集；最后，闭区间 \([a, b]\) （相应地开区间 \(]a, b[\) ）是 \([a, \rightarrow[\) 与 \(]\leftarrow, b]\) （相应地 \(]a, \rightarrow[\) 与 \(]\leftarrow, b[\) ）的交，因而是闭（相应地开）集。

因此 R 中的闭区间 \([-a, +a]\) （a \textgreater{} 0）都是 o 的邻域。我们来证明：当 a 取遍 \(R_{+}^{*}\) 时，它们组成 o 的一个基本邻域系。为此只需确立下述命题：

命题 3. 当 r 取遍 \textgreater0 的有理数组成的集时，R 中的区间 \(s_{r} = [-r, +r]\) 组成 0 的一个基本邻域系。

由第三章 § 3 第 4 目的命题 7，把诸区间

\(S_{r} \cap Q = [-r, +r]\) （ \(Q\) 中的区间）在 R 中取闭包，就得到 o 在 R 中的一个基本邻域系。只要证明了 \(S_{r}\) 是 \(S_{r} \cap Q\) 的闭包，证明即告完成。而 \(S_{r}\) 在 \(R\) 中是闭的，因此只需证明：若 \(x\) 是满足 \(-r < x < r\) 的实数，则 \(x\) 属于 \(S_{r} \cap Q\) 的闭包。区间 \(]-r, +r[\) 是 \(R\) 中的开集，因而对所有充分小的邻域 \(V\) （ \(o\) 在 \(R\) 中的），我们有 \(x + V \subset ]-r, +r[\) ；但 \(Q\) 在 \(R\) 中稠密，故存在有理数 \(r' \in x + V\) ，于是 \(-r < r' < r\) ，从而 \(r' \in S_{r} \cap Q\) 。

推论. 实直线的每一点都有可数的基本邻域系。

命题 4. 若 \((x, y)\) 是任一满足 \(x < y\) 的实数对，则存在有理数 \(r\) 使得 \(x < r < y\) 。

由于 \(\mathbf{Q}\) 在 \(\mathbf{R}\) 中稠密，只需证明 \(]x, y[\) 非空；由平移，可设 \(x = 0\) 且 \(y > 0\) 。而 \(\mathbf{R}\) 是豪斯多夫空间，因此由命题 3，存在有理数 \(r > 0\) 使得 \(y \notin [-r, +r]\) ，这就蕴含 \(0 < r < y\) 。

命题 5. 设 I 是 R 中任一区间。则 R 的拓扑在 I 上诱导的拓扑由 I 的开区间生成（其中 I 被视为由关系 \(x \leqslant y\) 全序化）。

I 的每个开区间是 R 的某个开区间在 I 上的迹。这对有界区间是显然的，而 I 的无界区间 \(]a, \rightarrow[\) 是 R 的无界区间 \(]a, \rightarrow[\) 的迹。因此我们可以限于 I = R 的情形；但在此情形结论由命题 3 得出，因为点 \(x \in R\) 的每个邻域都包含一个开区间 \(]x - a, x + a[\) 。

注. 若 A 是 R 的稠密子集，则 R 的拓扑由端点属于 A 的开区间生成。因为若

\[
] x - a, \quad x + a [
\]

是一个包含 \(x\) 的开区间，则存在两点 \(y, z\) （ \(A\) 的）使得 \(x - a < y < x\) 且 \(x < z < x + a\) ；于是 \(]y, z[\) 包含 \(x\) 且含于 \(]x - a, x + a[\) 。这一证明还表明，所考虑的诸区间组成 \(R\) 的拓扑的一个基（第一章 § 1 第 3 目）。特别地，若取 \(A = Q\) ，则我们看到 \(R\) 的拓扑具有可数基。

